\section{胶子—规范链顶点的重正化}
\label{sec:1g-link}
为使下面的讨论明确起见，我们假设图~\ref{fig:div-topo-g}中拓扑 (c) 图形里的规范链始于带算符 $F^{\mu \nu}(\xi_{1z})$ 的任意坐标 $\xi_{1z}$，终于不带附加算符的另一任意坐标 $\xi_{2z}$。记“裸”耦合常数为 $g_s$，“裸”的胶子场算符、Faddeev-Popov 鬼场与反鬼场分别为符号 $A$、$c$ 与 $\bar{c}$，与该拓扑相关的非阿贝尔场的广义 Ward 恒等式可推得为~\cite{Collins:1984xc}
\begin{align}\label{eq:gwi1}
\begin{split}
\langle -i \partial_\lambda^y A_d^\lambda(y)[\Phi(\{\xi_{2z},\xi_{1z}\})]_{a b} \,  F_{b}^{\mu \nu}(\xi_{1z}) \rangle\\
=\langle g_s \bar{c}_d (y) c_e (\xi_{2z}) [ t_e \Phi(\{\xi_{2z},\xi_{1z}\})]_{a b} \, F_{b}^{\mu \nu}(\xi_{1z}) \rangle ,
\end{split}
\end{align}
其中 $t$ 表示伴随表示的 SU(3) 生成元。
式~$(\ref{eq:gwi1})$ 的图像表示见图~\ref{fig:gwi1}，其中“1PR”表示单粒子可约图。图~\ref{fig:gwi1}左边的拓扑与图~\ref{fig:div-topo-g}中图 (c) 的相同，只是与外胶子动量缩并并以坐标空间表达。图~\ref{fig:gwi1}右边第一项的拓扑在时空中是非定域的，因而没有紫外发散。此外，在对 QCD 拉氏量以及规范链相关拓扑（图~\ref{fig:div-topo-g}中的 (a)、(b)）完成重正化之后，1PR 图的紫外发散会相消。也就是说，式~$(\ref{eq:gwi1})$ 的广义 Ward 恒等式保证了：若与外胶子动量缩并，图~\ref{fig:div-topo-g}中的拓扑 (c) 没有紫外发散。

%%%%%%%%%%%%%%%%%%%%%%%%%%%%%%%
\begin{figure}[htb]
	\begin{center}
		\includegraphics[width=0.55\textwidth]{images/Green-1-g.pdf}
		\caption{\label{fig:gwi1} 式~$(\ref{eq:gwi1})$ 中广义 Ward 恒等式的图像表示；虚线代表鬼场。}
	\end{center}
\end{figure}
%%%%%%%%%%%%%%%%%%%%%%%%%%%%%%%

为理解图~\ref{fig:div-topo-g}中拓扑 (c) 紫外发散的重正化，我们把该拓扑的图形记为 $\Gamma^{\lambda \mu \nu}(p,n)$，可称之为胶子—规范链顶点，其中 $p$ 与 $\lambda$ 分别是外胶子的动量和洛伦兹指标。洛伦兹对称性结合 $\mu$、$\nu$ 之间的反对称性使我们能做一般分解 $\Gamma^{\lambda \mu \nu}(p,n)=\sum_{i=1}^4 c_i \Pi_i^{\lambda \mu \nu}$，
\begin{align}
\begin{split}
\Pi_1^{\lambda \mu \nu}&=g^{\mu\lambda} p^\nu-g^{\nu\lambda}p^\mu,\, \Pi_2^{\lambda \mu \nu}=(p^\mu n^\nu-p^\nu n^\mu)n^\lambda,\\
\Pi_3^{\lambda \mu \nu}&=(p^\mu n^\nu-p^\nu n^\mu)p^\lambda, \,\Pi_4^{\lambda \mu \nu}=g^{\mu\lambda} n^\nu-g^{\nu\lambda}n^\mu.
\end{split}
\end{align}
如上所述，对紫外发散项有 $p_\lambda \Gamma^{\lambda \mu \nu}=0$，于是得到 $c_2\, p\cdot n + c_3\, p^2 + c_4=0$，进而
\begin{align}
\begin{split}
\Gamma^{\lambda \mu \nu}(p,n)=&c_1 \Pi_1^{\lambda \mu \nu} + c_2 (\Pi_2^{\lambda \mu \nu}-p\cdot n \Pi_4^{\lambda \mu \nu})\\
&+c_3 (\Pi_3^{\lambda \mu \nu}-p^2 \Pi_4^{\lambda \mu \nu}).
\end{split}
\end{align}
由于 $c_1$ 与 $c_2$ 的质量量纲为 $0$，紫外发散的定域性保证它们至多是对数发散；而 $c_3$ 因质量量纲为 $-1$ 而 UV 有限。唯一潜在线性紫外发散的系数 $c_4$ 被规范不变性消去。因此我们已经证明：图~\ref{fig:oneloop-g}中图 (b) 与 (c) 在一圈阶线性紫外发散的相消可以推广到所有阶，这使得准胶子算符的乘法可重正性成为可能。

在 $\alpha_s$ 的最低阶，有 $\Gamma^{\lambda \mu \nu}(p,n) \propto \Pi_1^{\lambda \mu \nu}$。如果要求 $\Gamma^{\lambda \mu \nu}(p,n)$ 在重正化下不与其他算符混合，就需要它的紫外发散在全阶都与 $\Pi_1^{\lambda \mu \nu}$ 成正比。幸运的是，事实确实如此。对于 $\mu$（或 $\nu$）沿 $z$ 方向的情形，或 $\mu$、$\nu$ 都不沿 $z$ 方向的情形，$c_2$ 项分别正比于 $\Pi_1^{\lambda \mu \nu}$ 或等于零。因此，$\Gamma^{\lambda \mu \nu}(p,n)$ 的各分量彼此完全不振荡混合，尽管两种不同选取需要两个不同的重正化常数。

综上，若为胶子—规范链顶点 $\Gamma^{\lambda \mu \nu}$ 选择 $F^{z \bar{\nu}}$ 或 $F^{\bar{\mu} \bar{\nu}}$，我们就可以通过乘以相应的重正化因子 $Z_{vg1}^{-1/2}$ 或 $Z_{vg2}^{-1/2}$ 分别消去该顶点的紫外发散。
